%% kapa: exam=1 round=37 year=2026 period=1 form=A subject=economics subject_ko=경제학원론 seq=38 number=78
%% source: Q-Net 감정평가사 기출문제 (cst00309) artlSeq=5258924 file=raw/1008/5258924/2026년 제37회 감정평가사 1차 1교시.pdf
%% conversion: kordoc 4.24.51 -> tools/gen1.py (auto); review=proofread-sonnet; flags=none
\begin{question}{78}
개방경제 균형국민소득 결정모형이 아래와 같다. 정부지출과 세금을 동시에 50씩 늘리는 경우 \uline{(ㄱ)균형소득의 변화량}과 \uline{(ㄴ)경상수지의 변화량}은? (단, C는 소비, Y는 국민소득, T는 세금, I는 투자, G는 정부지출, X는 수출, M은 수입이며, 재정정책 변화 전의 경제상태는 균형이었다.)
\begin{framedtext}
\begin{tabular}{@{}p{0.31\linewidth}@{\hspace{0.02\linewidth}}p{0.31\linewidth}@{\hspace{0.02\linewidth}}p{0.31\linewidth}@{}}
○ C = 2,000 + 0.8(Y - T) & ○ I = 900 & ○ G = 550 \\
○ T = 500 & ○ X = 500 & ○ M = 0.1(Y - T)
\end{tabular}
\end{framedtext}
\begin{choices}
  \item ㄱ: 50, ㄴ: 0
  \item ㄱ: 50, ㄴ: -5
  \item ㄱ: 100, ㄴ: -10
  \item ㄱ: 250, ㄴ: 0
  \item ㄱ: 250, ㄴ: -20
\end{choices}
\end{question}
\answer{1}
